% ============================================================
% 考题关键词（杂志页码 64–67）
% ============================================================
\colchapter{考题关键词}

\section*{“大思政课”}

中国共产党立志于中华民族千秋伟业，必须培养一代又一代拥护党的领导和我国社会主义制度、立志为中国特色社会主义事业奋斗终身的有用人才。思想政治理论课是落实立德树人根本任务的关键课程。习近平总书记曾深刻指出：“在大中小学循序渐进、螺旋上升地开设思政课非常必要，是培养一代又一代社会主义建设者和接班人的重要保障。”

当前，我国思政课建设整体上进入质量提升和内涵发展的新阶段。前不久，教育部等十部门印发《全面推进“大思政课”建设的工作方案》，旨在以“大思政课”建设为抓手，持续推动思政课和思想政治教育高质量发展。从思政课到“大思政课”，一字之别，体现了办好思政课的理念再更新、视野再开阔、格局再拓展。这就要求坚持开门办思政课，强化问题意识、突出实践导向，同时要推动思政小课堂与社会大课堂相结合，推动各类课程与思政课同向同行。

\section*{国家网络安全宣传周}

网络无远弗届，安全如影随形。9月5日至11日，2022年国家网络安全宣传周在全国范围内统一开展，主题为“网络安全为人民，网络安全靠人民”。

我国是名副其实的网络大国。截至2022年6月，我国网民规模已达10.51亿，人均每周上网时长为29.5个小时，形成了全球最为庞大、生机勃勃的数字社会。而网信事业越是繁荣发展，越不能忽视与日俱增的风险挑战。在今天，一个App，就能让手机里的信息如同“裸奔”；一条链接，就能让网络电信诈骗肆无忌惮……大到国家机密，小到个人隐私，网络安全防线稍有松懈，就有可能造成无法估量的损失。

安全是发展的前提，发展是安全的保障。只有处理好安全与发展的关系，做到协调一致、齐头并进，以安全保发展、以发展促安全，才能建久安之势、成长治之业。以国家网络安全宣传周为契机，让我们共同努力，一起筑牢“防火墙”，拧紧“安全阀”。

\clearpage
\section*{《羊了个羊》}

最近，小程序游戏《羊了个羊》突然爆火出圈，连续多日登上热搜。据悉，该游戏采用“消消乐”式玩法，玩家需将游戏页面上的众多图片按照三个一组的方式进行消除，直至页面再无图片。不少玩过的网友表示，这款游戏很“上头”，第一关极为简单，但第二关难度飙升，让人一度怀疑自己的智商，非要战胜不可。

《羊了个羊》不是第一款爆红的小程序游戏，更不会是最后一款。这种玩法简单、不需要专门客户端、没有内购项目，只让玩家看广告的游戏，被统称为“超休闲游戏”。其中达成“刷屏全网”成就的，国内近的就有去年的《合成大西瓜》，年代久远点的则是《旅行青蛙》。

比起创造一款截然不同的游戏，“超休闲游戏”的发力之处是瞄准玩家心理，知道他们喜欢什么、想要什么，他们愿意把时间花在哪里。《羊了个羊》的病毒式传播，正是因为给了我们一个机会，无须花费现实社交中的大量精力成本，亦能随时加入话题交流，而玩游戏只是一个借口。

\subsection*{巧命题}

最近，小程序游戏《羊了个羊》突然爆火出圈，连续多日登上热搜。单从话题性与传播度来看，《羊了个羊》无疑是成功的：无论是用奇低通关率激发用户胜负欲，还是属地捆绑撩拨用户“地域攀比心态”，它都把人的心理拿捏得死死的。可对于玩家来说，就算这类休闲游戏再“上头”，也要避免“贪杯”。比起游戏通关，更该打通的是“游戏沉迷”这道关，否则，我们只能沦为自己赌徒心态的“囚徒”。

对此，你有什么思考？请写一篇论述类文章加以阐明。

\subsection*{时材立意}

本命题材料介绍了《羊了个羊》成为爆款的原因——抓住了人的心理，同时告诫人们理性娱乐。就此，我们可得出如下立意：

\textbf{①控制欲望，适度游戏。}当我们在现实生活中面临诸多痛苦或者无聊时，有人便转向虚拟世界寻求快乐和成就感。超休闲游戏则精准拿捏了用户想逃避痛苦、追求快乐的心理诉求，持续进入游戏、接受任务、快速升级，不断提高正向反馈，再予以适当任务制造压力，而完成任务就可释放压力。这样的奖赏通路，正是上瘾行为的生理基础。如果我们不能节制自己的欲望，只会一路沉迷于此。

\textbf{②避免即时满足陷阱，建立延迟满足机制。}即时满足的快乐可以轻易获得，如打游戏、追剧、抽烟、酗酒……在现代社会，虚度光阴和获得满足有无数种方式。只不过，满足感的门槛一旦升高就很难降低，而门槛越高就越不容易感到满足，陷入空虚、无聊后又加大力度去探寻即时满足的快乐，从而造成恶性循环。只追寻即时快乐的最终结果只会让人愈加痛苦。我们只有甘愿为更有价值的长远结果而放弃部分即时满足的快乐，才能获得丰厚的回馈。

\textbf{③做人生这场游戏的“高质量玩家”，不负光阴。}醉心于超休闲游戏不会对人生有丝毫裨益，很多人沉溺其中，是因为没有找到真正喜欢的东西。生命可贵，可以做更多有意义的事，正青春的我们，当莫负好时光。新时代的舞台上处处刻印着青年奋斗的脚步，洋溢着青春火热的激情，当青年理想与历史使命同频共振，奏响的就是这个时代的最强音。

\clearpage
\section*{范文审阅：与其沉溺《羊了个羊》，不如奋起“学了个学”}
\auth{文/山东省青岛市胶南一中 余依然（得分：57分）}

最近一款小游戏火爆朋友圈，不少网友很“上头”，直接沉溺其中无法自拔，追寻着虚拟的满足感。据传，因为玩家太多，服务器几度崩溃；为了闯过第二关，好多网友通宵达旦……对于这种现象，我认为，与其沉溺于无意义的虚拟世界，享受即时满足感，不如将光阴用在更有意义的地方。

诚然，这款小游戏能让你产生极大的兴趣。《羊了个羊》共两关，第一关玩家可轻松通过，让人拥有一种成就感；第二关则颇具难度，让人产生挑战的欲望，以收获更大的“荣誉感”——跟朋友交谈“闯关心得”，在熟人面前炫耀“通关秘诀”，极低的通关率又佐证了你的高智商……于是，纵然被虐千万遍，却依旧“初心不改”屡败屡战，游戏占据了你的大半空闲时光。

我们总是称赞执着和专注，巴菲特也曾将自己的成功归结于“专注”。可是，亲爱的朋友，你的专注是否用错了地方？沉溺于虚拟世界，沉溺于即时快感，沉溺于游戏，这些做法，是不是值得反思？庄子有言：“用志不分，乃凝于神。”周恩来总理曾说：“宁精勿杂，宁专勿多。”的确，业余时间适当娱乐有益于调节自我、放松心情，但你也要明白过犹不及、玩物丧志的道理。把专注用错了地方，就可能迷失方向，甚至走向歧途。

联系现实生活，我常常为一些不理智的行为而叹息：有学生因沉溺游戏不能自拔而荒废学业；有“打工人”为游戏通关而陷入网络骗局，遭受财产损失……与其沉溺《羊了个羊》，不如奋起“学了个学”。争分夺秒、努力学习，奋勇争先、为国奉献的青年们，才是一代国之脊梁！北京大学的90后学生宋玺，专注于有意义的事业，不爱红装爱武装，全力护航亚丁湾；云南80后干部李忠凯，专注在脱贫攻坚主战场，为祖国而忙碌，为人民而奔波，哪怕青丝变白发；河北保定支教青年奔赴大漠，专注地播种希望，让青春芳华有了绽放的意义……

一代人有一代人的使命，一代人有一代人的担当。作为当代青年，我们不可沉溺虚拟世界，须以青春之我、奋斗之我，积极投身于民族复兴的大业，燃青春梦想，展青春激情，成靓丽人生。

\score{得分点：巧用关联词，旗帜鲜明亮出中心论点。}
\score{得分点：联系现实，观点深刻。}
\score{得分点：摆事实讲道理，联系该款游戏的特点进行分析，说理充分。}
\score{得分点：引用论证彰显文采；连续反问突出思考。}
\score{得分点：引用当代杰出人物的事例，时代气息浓重，与上文形成明显对比，凸显人生价值。}

\subsection*{时评链接：沉迷《羊了个羊》，当心落入“虚拟狼圈”}

网络上这样形容《羊了个羊》的火爆程度：“月亮不睡我不睡，势必要过第二关。”在此，还是要告诫网友，游戏只是为生活调味，真正的“羊”，还得回归现实世界，与身边朋友一起“入圈”，做更加有意义的事，遇见更加美好的风景。

为一场没有实际意义的游戏熬夜，到头来还是一场空、一场虚梦。凡事皆有度，过犹不及。适当的游戏时间可以消解疲劳，能够给现实的世界带来快乐和激情。但是，若以疯狂的商业吹捧，裹挟起“全民皆羊”的不良之风，那么就打错了算盘。对此，公众要有清晰的认知和独立的思考——任何的爆款游戏，都无法替代现实的情感与体知，想要慰藉和欢喜，还得从现实的奋斗中一点一滴获取。沉迷《羊了个羊》不能自拔，只会浪费少年时，空悲切，不值当。

（摘自大江网，有删改）

\subsection*{名师态度}

本文按照“引、议、联、结”的结构模式娓娓道来，给人以深刻的启迪。文章开篇先介绍《羊了个羊》火爆朋友圈的热点现象，再由此剖析背后的原因，并合理过渡到分析沉溺游戏的危害，指出当代青年应该追寻有价值的事情。特别值得称道的是，本文作者注意锤炼语言，引用论证凸显文采，排比论证彰显深度，对偶句式凸显匠心。

——山东省青岛市胶南一中 徐增才
